\newcommand{\BalReach}{\mathcal{B}_{(4)}}
\newcommand{\autPlus}{\aut_{\sigma}}
We develop algorithms for deciding consistency (resp., positive or negative consistency). 
We discuss separately well-matched rules (Section~\ref{s:well-matched}) and the wider class of matching-preserving rules (Section~\ref{s:matching-preserving}).
In both cases, we consider first rules without variables, and then the extension to rules with variables.
Although the asymptotic complexity is the same in both cases, the algorithms for well-matched rules are conceptually simpler and admit more efficient implementations.
Before studying particular types of rules, we make a general observation that reduces consistency checking to single-step rewriting.

\Paragraph{Reduction to single-step rewriting}
 Observe that an \SRSV $\rew$ is consistent with the language of $\aut$ if and only if for every rule $\lhs \to \rhs$ and all words $x,y$, we have
$x \lhs y \in \lang(\aut) \iff x \rhs y \in \lang(\aut)$. Indeed, if this property holds for every rule, then 
 consistency under multi-step rewriting follows by transitivity.
 This holds due to the assumption that the $\aut$ and $\rew$ are over the same alphabet~\cite{ecai25}.
%by transitivity of the equivalence 
%if $w \rewrites_{\rew} w'$, then $w \in \lang \implies w' \in \lang$.
Analogously, for positive and negative consistency it suffices to verify the property for each individual rewriting step.

\subsection{Consistency of well-matched rules}
\label{s:well-matched}

\Paragraph{Ground well-matched rules} Consider a well-matched rule $\lhs \to \rhs$ without variables, i.e., a standard string rewriting rule. 
We extend the alphabet $\vAlph$ with a fresh local letter $\#$ and construct two \DVPA over the extended alphabet.
Both \DVPA extend $\aut$ and differ only in their behavior on the letter $\#$; in $\aut_1$ for any state $q$, 
the transition over $\#$ leads to the state $q'$ such that
$\aut$ moves from $q$ to $q'$ over $\lhs$. 
Since $\lhs$ is well-matched, the run over $\lhs$ does not depend on the stack content and it does not modify the stack content. 
Thus, the effect of reading $\lhs$ in the \DVPA $\aut$ is equivalent to a transition over a single local letter. 
The \DVPA $\aut_2$ is defined analogously, but $\#$ simulates the word $\rhs$ instead. Again, word $\rhs$ is well-matched. 
Therefore, to check consistency of  $\lhs \to \rhs$  with $\lang(\aut)$ it suffices to check whether there is a word $w$ such that exactly one of \DVPA $\aut_1, \aut_2$ accepts it, 
which can be done by a single reachability check on the product \DVPA $\aut_1 \times \aut_2$, a \DVPA simulating two copies of $\aut$ in parallel.
The visibility condition ensures that a single stack over the product of stack alphabets is sufficient to simulate both \DVPA.
Checking positive (resp., negative) consistency amounts to checking whether $\lang(\aut_1) \setminus \lang(\aut_2)$  (resp., $\lang(\aut_2) \setminus \lang(\aut_1)$ for negative consistency) is non-empty, 
which again amounts to a single reachability check. 

\begin{proposition}
Consistency (resp.,  positive, negative consistency) of a ground well-matched rule with the language of a given \DVPA $\aut$ can be decided in polynomial time via a single 
reachability check on a \DVPA obtained from the product $\aut \times \aut$. 
\end{proposition}


Reachability in pushdown automata can be checked in $\bigO(|\aut|^3)$, with a slightly improved bound $\bigO(|\aut|^3/\log(|\aut|))$ for \emph{recursive state machines}~\cite{DBLP:conf/popl/Chaudhuri08}
both of which fall into $\widetilde{\bigO}(|\aut|^3)$.
Due to the lack of better upper bounds~\cite{DBLP:conf/popl/Chaudhuri08,DBLP:conf/lics/BalasubramanianCM25}, the complexity of reachability in pushdown systems is often referred to as the ``cubic bottleneck.''
Consequently, reachability in the product $\aut \times \aut$ amounts to $\widetilde{\bigO}(|\aut|^6)$. 
However, checking consistency of an SRS with the language of a \DFA can be performed locally on the structure of the \DFA without checking reachability~\cite{ecai25}. 
In contrast, the following example illustrates that for pushdown systems, checking consistency must involve reachability.

\begin{example}
Consider a \DVPA $\autB$ depicted in Figure~\ref{fig:nonsyntactic}, which is over the pushdown alphabet $(\set{c},\set{l_1, l_2}, \set{r})$ such that local letters $l_1, l_2$ induce self-loops on 
all states except for $q_0$. Its stack alphabet is a singleton $\set{a}$.
In the state $q_0$ the \DVPA moves over $l_1$ to $q_1$ and over $l_2$ to $q_2$. The automaton behaves differently in states $q_1, q_2$ if the stack is non-empty, while these states are indistinguishable 
if the stack is empty. 
Finally, $q_0$ is reachable only with non-empty stack, i.e., $q_0$ is reachable, but the configuration $(q_0, \bot)$ is not.
The rule $l_1 \to l_2$ is consistent with $\lang(\autB)$, but establishing this requires reachability analysis. 
\begin{figure}[h!]
\centering
\begin{tikzpicture}[
    shorten >=1pt, 
    node distance=3.0cm, 
    on grid, 
    auto,
    >={Stealth[bend]}
]

  % State Definitions
  \node[state,initial] (q0) {$q_I$};
  \node[state] (q1) [right=of q0] {$q_0$};
  \node[state] (q2U) [above right=of q1] {$q_1$};
  \node[state] (q2D) [below right=of q1] {$q_2$};
  \node[state] (sink) [below right=of q2U] {$q_R$};
  \node[state,accepting] (acc) [right=of q2U] {$q_F$};

  % Transition Edges
  \path[->]
    (q0) edge node [above] {$c, push(a)$} (q1)
    (q1) edge node [above] {$l_1$} (q2U)
    (q1)  edge node [above] {$l_2$} (q2D)
    (q2U) edge node [right] {$r, pop(a)$} (sink)
    (q2U) edge node [above] {$r, pop(\bot)$} (acc)
    (q2D)  edge [bend right=20] node [above left=0.6cm] {$r, pop(a)$} (sink) 
    (q2D)  edge [bend left=20] node [below right=0.6cm] {$r, pop(\bot)$} (sink);
    

    %[bend right=45]
 %   (q1)  edge [bend right=45] node [above] {$\Sigma\setminus (Y \cup \set{x})$} (qn)
    % The outgoing edge xi from q1
%    (q1)  edge [nodes={at={(1, -0.5)}}] node {$x$} ++(1.5, 0);
%    (q1)  edge node [above] {$x$} (qF)
%    (q1) edge node [left] {$Y$} (sink)
%    (q2) edge node [left] {$Y$} (sink)
%    (qn) edge node [above] {$Y$} (sink)
%    (qF) edge node [left] {$Y$} (sink)
    ;

 \draw[loop right,->] (acc)  to node[right] {$\vAlph$} (acc);
 \draw[loop right,->] (sink) to node[right] {$\vAlph$} (sink);

\end{tikzpicture}
\caption{The automaton $\autB$. All transitions other than depicted are self-loops}
\label{fig:nonsyntactic}
\end{figure}
\end{example} 


\Paragraph{Non-ground well-matched rules}
\newcommand{\autCon}{\aut_{\textrm{con}}}
\newcommand{\DeltaRule}{\Delta[\lhs \to \rhs]}
Consider a well-matched rule $\lhs \to \rhs$ with variables. 
The general approach, as in the ground case, is to reduce consistency checking to reachability in a product \VPA.
We construct a \VPA $\autCon$ that simulates two copies of $\aut$ over the alphabet extended with a fresh local letter $\#$.
The \VPA $\autCon$ is built on the product $\aut \times \aut$, whose stack track contents of both copies of $\aut$.
The fresh local letter $\#$ simulates the word $\lhs$ in the first copy and the word $\rhs$ in the second copy.
However, there are two key differences.
First, since $\lhs$ and $\rhs$ can be instantiated to (infinitely many) different words, transitions over $\#$ are non-deterministic.
Second, the transitions in each copy of $\aut$ cannot be defined independently; they must be defined jointly on the product to ensure that both sides of the rule $\lhs \to \rhs$ are instantiated to the same word. 
For example, for a rule $c X r \to X$, there may be a transition in $\aut$ from $q$ to $q_1$ over $c w_1 r$ and from $q$ to $q_2$ over $w_2$, but we must ensure that $w_1 = w_2$.
To this end, we compute a set $\DeltaRule$ of triples $(q, s_1, s_2)$ such that there are words $\alpha, \beta$ satisfying:
(1)~$s_1 \neq s_2$, 
(3)~$(q,u) \trans{\aut}{\alpha} (s_1,u)$ and $(q,u) \trans{\aut}{\beta} (s_2,u)$ for every stack content $u$, and
(4)~$\alpha \to \beta$ is an instance  of $\lhs \to \rhs$ obtained by substituting well-matched words for variables.
Using $\DeltaRule$, we define a \VPA $\autCon$ over the pushdown alphabet $\vAlph$ extended with a local letter $\#$ as an extension of $\aut \times \aut$ such that for a state $(q_1, q_2)$ of $\autCon$:
(a)~if $q_1 \neq q_2$, then $\autCon$ has no transition  over $\#$, and
(b)~if $q_1 = q_2$, then for every $(q_1, s_1, s_2) \in \DeltaRule$, 
the tuple $\langle (q_1, q_1), \#, (s_1, s_2) \rangle$ is a transition in $\autCon$. 
Intuitively, the automaton only applies the rule when both copies are in the same state, and hence words with more than one $\#$ letter are rejected.
Finally, the automaton accepts if in the reached state $(q_1, q_2)$, exactly one of $q_1, q_2$ is accepting in $\aut$. 
Observe that the \VPA $\autCon$ accepts exactly words $w$ that are witnesses of inconsistency of $\lhs \to \rhs$ with $\lang(\aut)$, i.e., 
any accepted word $w$ contains exactly one occurrence of $\#$. 
Furthermore, there is a ground instance $\alpha \to \beta$ of the rule $\lhs \to \rhs$ obtained by substituting well-matched word for variables, and 
by substituting $\#$ with $\alpha$ (resp., $\beta$) yields $w_1$ (resp., $w_2$) satisfying 
(i)~$w_1$ can be rewritten on one step to $w_2$ with the rule $\alpha \to \beta$, and 
(ii)~exactly one of $w_1, w_2$ belongs to $\lang(\aut)$.
By modifying the acceptance condition, we can also decide positive and negative consistency. 

In the following, we focus on computing $\DeltaRule$ for rules with variables. 
It suffices to compute how well-matched words transform pairs of states of the \DVPA.

\Paragraph{Computing $\DeltaRule$}
\newcommand{\lhsInst}{\bar{\alpha}}
\newcommand{\rhsInst}{\bar{\beta}}
We compute a quaternary relation $\BalReach \subseteq Q^4$ such that $(q_1, q_2, s_1, s_2) \in \BalReach$ if and only if there is a well-matched word $w$ such that
$q_1 \trans{\aut}{w} s_1$ and $q_2 \trans{\aut}{w} s_2$. 
It is computed iteratively in a similar way to reachability in pushdown automata, but applied to the product automaton $\aut \times \aut$.
 %we compute the closure with respect to 
%enclosing the word with a call and a return letter, joining a local letter and concatenation computed simultaneously.
\begin{restatable}{lemma}{BalancedReachability}
Given a \DVPA $\aut$, the relation $\BalReach$ can be computed in polynomial time in the size of $\aut$.
\end{restatable}
We extend the alphabet $\vAlph$ to $\vAlph_+$ by adding a pair of fresh local letters $x^L$ and $x^R$ for every variable $X$ occurring in $\lhs \to \rhs$.
Next, we substitute each variable $X$ with its corresponding local letters: $x^L$ for $X$ in $\lhs$ and $x^R$ for $X$ in $\rhs$, obtaining words $\lhsInst$ and $\rhsInst$ over the extended alphabet. 
Then, we iterate over all possible assignments from variables of $\lhs \to \rhs$ to $\BalReach$. 
For a given assignment $\sigma$, we extend the \DVPA $\aut$ to a \VPA $\autPlus$ working over the extended alphabet, where transitions over new letters are defined as follows.
For every variable $X$, if $\sigma(X) = (q_1, q_2, s_1, s_2)$, then
$\autPlus$ has a transition over $x^L$ from $q_1$ to $s_1$ and over $x^R$ from $q_2$ to $s_2$; there are no other transitions over $x^L$ or $x^R$.
Finally, %we evaluate $\lhs$ and $\rhs$ by evaluating $\lhsInst$ and $\beta^+$ in $\autPlus$.
for every triple $(q,s_1, s_2)$, if $(q,\bot) \trans{\autPlus}{\lhsInst} (s_1,\bot)$ and $(q,\bot) \trans{\autPlus}{\rhsInst} (s_2,\bot)$, we add the triple $(q,s_1, s_2)$  to $\DeltaRule$.
Note that the stack content $\bot$ does not influence reachability since $\lhsInst, \rhsInst$ are well-matched.


\begin{lemma}
Deciding membership in $\DeltaRule$ over well-matched rules is in \NP.
Under a fixed bound on the number of variables per rule, it is decidable in polynomial time.
\end{lemma}

In consequence, checking inconsistency can be done in \NP, which gives \coNP test for consistency:

\begin{theorem}
Consistency (resp., positive, negative consistency) of a well-matched rule with the language of a given \DVPA $\aut$ can be decided in $\coNP$.
Under a fixed bound on the number of variables per rule, it is decidable in polynomial time.
\end{theorem}


\Paragraph{Complexity of checking (positive) consistency}
The proposed algorithm checks all assignments of transitions for variables, which is exponential in the number of variables; it works in polynomial time if the number of variables is bounded. 
Significantly improving upon this complexity is unlikely as the membership problem for $\DeltaRule$ is \NP-complete in general, and 
deciding positive (resp. negative) consistency is \coNP-complete.
The hardness result holds already for \DFA, which can be considered as \DVPA over only local letters.


\begin{restatable}{theorem}{NPHardness}
(1)~The problem of deciding, given a rewrite rule $\lhs \to \rhs$, a DFA $\aut$ and its states $q, s_1, s_2$, whether $(q, s_1, s_2)$ belongs to $\DeltaRule$, is \NP-complete. 
(2)~The positive (resp. negative) consistency problem over well-matched rules against a regular language given by a \DFA is \coNP-complete.
\end{restatable}
\begin{proof}[Proof's sketch]
We show thus by reduction from the 3-SAT problem. Given a formula $\varphi$, we rename variables to obtain $\varphi^*$, in which  each variable occurs at most once. 
Then, $\aut$ moves over some instance of $\lhs$ from $q$ to $s_1$ if the variable substitution of $\lhs$ encodes a truth assignment, where all copies of the same variable have the same value.
Next, $\aut$ moves over some instance of $\rhs$ from $q$ to $s_2$  if  the variable substitution of $\rhs$ encodes a truth assignment satisfying the formula $\varphi$.
Thus, $(q, s_1, s_2)$ belongs to $\DeltaRule$ if and only if $\varphi$ is satisfiable.
Furthermore, we can reduce unsatisfiablity of $\varphi$ to positive (resp. negative) consistency. 
We ensure that whenever a truth assignment is consistent among copies of the same propositional variable, it does
not satisfy $\varphi^*$. Thus, $\varphi$ is unsatisfiable.
\end{proof}

The \coNP-hardness result for positive (resp. negative) consistency does not extend to consistency, and hence 
the complexity lower bound for checking consistency  of well-matched rules remains open.

\subsection{Consistency of matching-preserving rules}
\label{s:matching-preserving}

We turn to matching-preserving rules that are not well-matched.
The overall approach is similar to the well-matched case; we extend the alphabet with auxiliary letters and consider the product $\aut \times \aut$.
The key difference is that in well-matched rules, both sides are stack-independent and stack-neutral, whereas matching-preserving rules may depend on and modify the stack.
As a result, instead of local letters, we extend the alphabet with call and return letters. 

\Paragraph{Ground matching-preserving rules}
\newcommand{\autExt}{\aut_E}
\newcommand{\return}{\overline{\mathbf{r}}}
\newcommand{\call}{\overline{\mathbf{c}}}
Consider a matching-preserving rule $\lhs \to \rhs$ without variables that is not well-matched.
Since the rule is matching-preserving, $\lhs$ and $\rhs$ have the same number $m$ of unmatched calls and the same number $k$ of unmatched returns; 
since it is not well-matched, there are unmatched letters (at least one of $k, m$ is non-zero); assume that $k>0$ (the case $k=0$ and $m>0$ is similar). 
Furthermore,  $\lhs$ and $\rhs$ can be partitioned into well-matched words and unmatched calls and returns, where
unmatched returns precede unmatched calls.
It follows that  
\begin{align*}
&&\lhs &= v_0^L \cdot r_1^L \cdot v_1^L \ldots r_k^L \cdot v_k^L \cdot c_1^L \cdot v_{k+1}^L \cdot c_2^L \ldots c_m^L \cdot v_{k+m}^L \\
&&\rhs &= v_0^R \cdot r_1^R \cdot v_1^R \ldots r_k^R \cdot v_k^R \cdot c_1^R \cdot v_{k+1}^R \cdot c_2^R \ldots c_m^R \cdot v_{k+m}^R
\end{align*}
where $v_i^L, v_i^R$ are well-matched words, $r_i^L, r_i^R$ are return letters and $c_i^L,c_i^R$ are call letters. 
As in the well-matched case, let $\autExt$ be an extension of the product automaton $\aut \times \aut$, i.e., 
its set of states is $Q \times Q$ and its stack alphabet is $\stackAlphabet \times \stackAlphabet$.
The automaton $\autExt$ works over an extended alphabet, where 
the additional letters are $k$ fresh return letters
$\return_1, \ldots, \return_k$
and $m$ fresh call letters $\call_1, \ldots, \call_m$. 
The transition over $\return_1$ simultaneously  simulates reading $v_0^L r_1^L v_1^L$  on the first component of $\autExt$ and $v_0^R r_1^R v_1^R$ on the second.
In doing so, it pops one symbol from the product stack, which encodes a pair of stack symbols from the two copies of $\aut$.
Formally, $\delta_E((q_1, q_2),\return_i, (b_1, b_2)) = (s_1, s_2)$ if $(q_1,u b_1) \trans{\aut}{v_0^L r_1^L v_1^L} (s_1, u)$ and 
$(q_2,u b_2) \trans{\aut}{v_0^R r_1^R v_1^R} (s_2, u)$ for every stack content $u$. 
For the remaining return letters $\return_i$, where $i \in \set{2, \ldots,k}$,
the transitions on $\return_i$ simulate reading  $r_i^L v_i^L$  on the first component and $r_i^R v_i^R$ on the second component.
For each new call letter $\call_i$, the transition over $\call_i$ simultaneously  simulates reading $c_i^L v_{k+i}^L$ on the first component and $c_i^R v_{i+k}^R$ on the second.
Finally, a state $(q_1, q_2)$ of $\autExt$ is accepting if exactly one of $q_1, q_2$ is accepting in $\aut$.

The following condition holds:
\begin{claim}
A  word  $w_1 \return_1 \ldots \return_k \call_1 \cdots \call_m w_2$, where $w_1, w_2$ are over the original alphabet $\vAlph$, is accepted by $\autExt$  
if and only if $\aut$ accepts exactly one of the words $w_1 \lhs w_2, w_1 \rhs w_2$, which holds precisely when $\lhs \to \rhs$ is inconsistent with $\lang(\aut)$.
\end{claim}
Therefore, the rule $\lhs \to \rhs$ is consistent with $\lang(\aut)$ if and only if
the language $\lang(\autExt) \cap \lang(\vAlph^* \return_1 \ldots \return_k \call_1 \cdots \call_m \vAlph^*)$ is empty.

By changing the accepting states of $\autExt$ as discussed in the well-matched case, we can decide positive and negative consistency 
using emptiness of $\lang(\autExt) \cap \lang(\vAlph^* \return_1 \ldots \return_k \call_1 \cdots \call_m \vAlph^*)$.
In consequence, we have:

\begin{proposition}
Consistency (resp.,  positive, negative consistency) of a ground matching-preserving rule with the language of a given \DVPA $\aut$ can be decided in polynomial time via a single reachability check on a \DVPA obtained from the product $\aut \times \aut \times \autB$, where $\autB$ is a \DFA depending only on the rule.
\end{proposition}

Now, we can move to rules with variables. We only briefly comment on the extension as it is similar as for well-matched rules.

\Paragraph{Matching-preserving rules with variables} 
Variables in matching-preserving rules can be eliminated as in the case of well-matched rules. 
Consider a matching-preserving rule $\lhs \to \rhs$.
It suffices to extend the alphabet with local letters $x^L, x^R$ for each variable and then consider a ground rule
obtained from  $\lhs \to \rhs$ by replacing the occurrence of every variable $X$ with $x^L$ in $\lhs$ and with $x^R$ in $\rhs$ whenever appears.
Finally, as described in Section~\ref{s:well-matched} for the well-matched case, we iterate over all possible assignments from variables of $\lhs \to \rhs$ to $\BalReach$. 
In consequence, we have the following:

\begin{theorem}
Consistency (resp., positive, negative consistency) of a matching-preserving rule with the language of a given \DVPA $\aut$ can be decided in $\coNP$.
Under a fixed bound on the number of variables per rule, it is decidable in polynomial time.
\end{theorem}




%%% IMPROVMENTS

%The complexity discussion is split across two locations in the well-matched subsection. The proposition on ground rules states polynomial-time decidability, then the cubic bottleneck discussion follows, then the example, then the non-ground case introduces NP membership, and finally the hardness results appear at the end. This means the reader encounters upper bounds, then a digression, then more upper bounds, then lower bounds — all interleaved with algorithmic content. Consider consolidating the complexity discussion: state all upper bounds first (ground: polynomial, non-ground: coNP / polynomial with bounded variables), then present the hardness results together at the end of the subsection. This would give the reader a cleaner picture.

%The example is well-motivated but interrupts the algorithmic narrative. It appears between the ground and non-ground cases, which is a natural break point. However, it makes the point that reachability is necessary, which is a lower-bound-style observation. Placing it after the complexity results — rather than between the two algorithmic cases — would keep the algorithmic development uninterrupted and group all complexity-related discussion together.

%The matching-preserving subsection feels noticeably shorter and less developed than the well-matched one. The ground case is explained in reasonable detail, but the treatment of variables is compressed into two sentences that simply say "do the same as before." While this is accurate, it leaves the reader unsure whether there are genuinely no new complications. If the argument is indeed identical, a single explicit sentence confirming this — "The construction for variables in matching-preserving rules introduces no new difficulties beyond those already handled in the well-matched case" — would be more reassuring than a terse forward reference. If there are subtle differences, they should be flagged.

%The two subsections are not fully parallel in structure. The well-matched subsection has: ground rules → proposition → complexity discussion → example → non-ground rules → \DeltaRule computation → lemma → theorem → hardness discussion. The matching-preserving subsection has: ground rules → proposition → variables in two sentences → theorem. The asymmetry is partly justified by the fact that the variable elimination is identical, but the reader may feel that something is missing. Either making the structure more explicitly parallel, or acknowledging the asymmetry and explaining why it is justified, would improve the overall coherence.

% The open problem at the end of the well-matched subsection is a reasonable closing remark, but it lands somewhat flatly. Since it concerns only the lower bound for full consistency (not positive/negative), it is a precise and meaningful observation. Consider framing it slightly more prominently — for instance, as a named open question — so it does not read as an afterthought.